\documentclass[11pt,a4paper]{article}
\usepackage[utf8]{inputenc}
\usepackage{amsmath,amssymb,amsthm}
\usepackage{listings}
\usepackage[margin=2.5cm]{geometry}
\usepackage{xcolor}

\lstset{
  language=C++,
  basicstyle=\ttfamily\small,
  keywordstyle=\color{blue}\bfseries,
  commentstyle=\color{green!50!black},
  stringstyle=\color{red!70!black},
  numbers=left,
  numberstyle=\tiny\color{gray},
  breaklines=true,
  frame=single,
  tabsize=4,
  showstringspaces=false
}

\title{IOI 1991 -- Problem 1: Island}
\author{Solution}
\date{}

\begin{document}
\maketitle

\section{Problem Statement}

Given an $R \times C$ grid where each cell is either land (\texttt{1}) or
water (\texttt{0}), count the number of islands. An island is a maximal
connected component of land cells under 4-directional adjacency (up, down,
left, right).

\section{Solution}

This is the classic connected-components-on-a-grid problem, solved by flood
fill.

\subsection{Algorithm}

\begin{enumerate}
  \item Iterate through every cell $(i, j)$ in the grid.
  \item When an unvisited land cell is found, increment the island counter
    and perform a BFS (or DFS) to mark all reachable land cells as visited.
  \item The final counter equals the number of islands.
\end{enumerate}

\subsection{Correctness}

Each BFS/DFS from an unvisited land cell discovers exactly one maximal
connected component, since it visits all cells reachable via 4-adjacency
from the starting cell. Two distinct BFS calls never visit the same cell,
so each component is counted exactly once.

\subsection{BFS vs.\ DFS}

BFS (using a queue) avoids the risk of stack overflow on large grids, which
can occur with recursive DFS. We use BFS here.

\section{Complexity Analysis}

\begin{itemize}
  \item \textbf{Time:} $O(RC)$. Each cell is enqueued and processed at most
    once.
  \item \textbf{Space:} $O(RC)$ for the visited array. The BFS queue holds
    at most $O(\min(R, C))$ elements for typical grid shapes, though
    worst-case it can hold $O(RC)$ elements.
\end{itemize}

\section{Implementation}

\begin{lstlisting}
#include <iostream>
#include <vector>
#include <queue>
using namespace std;

int main() {
    int R, C;
    cin >> R >> C;

    vector<vector<int>> grid(R, vector<int>(C));
    vector<vector<bool>> visited(R, vector<bool>(C, false));

    for (int i = 0; i < R; i++)
        for (int j = 0; j < C; j++)
            cin >> grid[i][j];

    const int dx[] = {0, 0, 1, -1};
    const int dy[] = {1, -1, 0, 0};

    int islands = 0;

    for (int i = 0; i < R; i++) {
        for (int j = 0; j < C; j++) {
            if (grid[i][j] == 1 && !visited[i][j]) {
                islands++;
                queue<pair<int,int>> q;
                q.push({i, j});
                visited[i][j] = true;

                while (!q.empty()) {
                    auto [r, c] = q.front();
                    q.pop();
                    for (int d = 0; d < 4; d++) {
                        int nr = r + dx[d];
                        int nc = c + dy[d];
                        if (nr >= 0 && nr < R && nc >= 0 && nc < C
                            && !visited[nr][nc]
                            && grid[nr][nc] == 1) {
                            visited[nr][nc] = true;
                            q.push({nr, nc});
                        }
                    }
                }
            }
        }
    }

    cout << islands << endl;
    return 0;
}
\end{lstlisting}

\section{Example}

\begin{verbatim}
Input (5 x 5):
1 1 0 0 0
1 1 0 0 1
0 0 0 1 1
0 0 0 0 0
1 0 1 0 1

Output: 5
\end{verbatim}

The five islands are: $\{(0,0),(0,1),(1,0),(1,1)\}$,
$\{(1,4),(2,3),(2,4)\}$, $\{(4,0)\}$, $\{(4,2)\}$, and $\{(4,4)\}$.

\end{document}
